% =========================================================
%  Research-RAG: A Hybrid Retrieval-Augmented Generation
%  System for Multi-Paper Academic Question Answering
%  IEEE Conference Paper — IEEEtran class
% =========================================================
\documentclass[conference]{IEEEtran}

% ---- Packages ------------------------------------------
\usepackage[T1]{fontenc}
\usepackage[utf8]{inputenc}
\usepackage{cite}
\usepackage{amsmath,amssymb,amsfonts}
\usepackage{algorithmic}
\usepackage{graphicx}
\usepackage{textcomp}
\usepackage{xcolor}
\usepackage{hyperref}
\usepackage{booktabs}
\usepackage{multirow}
\usepackage{array}
\usepackage{url}
\usepackage{algorithm}
\usepackage{algpseudocode}

\hypersetup{
  colorlinks=true,
  linkcolor=blue,
  citecolor=blue,
  urlcolor=blue
}

% =========================================================
\begin{document}

\title{Research-RAG: A Hybrid Knowledge-Graph and Dense--Sparse\\
Retrieval-Augmented Generation Framework for\\
Multi-Paper Academic Question Answering}

\author{
  \IEEEauthorblockN{[Author Name]}
  \IEEEauthorblockA{
    \textit{Department of Computer Science and Engineering}\\
    \textit{[Institution Name]}\\
    [City, Country]\\
    [email@institution.edu]
  }
}

\maketitle

% =========================================================
\begin{abstract}
Retrieval-Augmented Generation (RAG) has emerged as a powerful paradigm
for grounding large language model (LLM) responses in verifiable source
material. However, conventional dense-retrieval RAG pipelines treat
documents as flat bags of chunks, discarding the rich relational structure
that links methods, datasets, metrics, and empirical results within and
across research papers. We present \textbf{Research-RAG}, a hybrid RAG
system that combines (i) a structured knowledge graph built by LLM-driven
entity extraction from PDF documents, (ii) FAISS-based dense semantic
search over document chunks, and (iii) BM25 sparse retrieval, fused via a
weighted hybrid score. The system supports multi-paper ingestion, produces
grounded structured JSON answers, and includes a built-in hallucination
detection module.

We evaluate Research-RAG against a na{\"i}ve dense-only baseline
(Simple~RAG) and a pure graph-traversal approach (Graph~RAG) on 12
questions across 4 research papers, scored on two axes by an independent
LLM judge: Answer Faithfulness and Completeness.
The results show that each retrieval modality has a distinct strength:
Simple~RAG achieves the highest overall Combined score (0.8458) due to
consistently complete passage-level evidence; Graph~RAG excels on
structured-fact questions (Q12: Combined~=~1.00); and Hybrid~RAG
achieves the strongest Reasoning scores (Combined~=~0.850), outperforming
Graph~RAG by~+0.2375 on that subset.
The experiments reveal a completeness shortfall in pure graph retrieval
for prose-heavy reasoning queries, confirming the complementary value of
fusing all three modalities.
\end{abstract}

\begin{IEEEkeywords}
Retrieval-Augmented Generation, Knowledge Graphs, Dense Retrieval,
BM25, Multi-Paper Question Answering, LLM, FAISS, NetworkX
\end{IEEEkeywords}

% =========================================================
\section{Introduction}
\label{sec:intro}

Large language models (LLMs) such as GPT-4 \cite{openai2023gpt4} and
Gemini \cite{team2023gemini} have demonstrated remarkable capabilities in
language understanding and generation. However, they are prone to
hallucinations---confidently producing statements that are factually
incorrect or unsupported by source material \cite{ji2023survey}. This
limitation is particularly pronounced in knowledge-intensive domains such
as academic research, where precision, citation accuracy, and cross-paper
comparison are essential.

Retrieval-Augmented Generation (RAG) \cite{lewis2020rag} addresses this by
conditioning generation on retrieved external passages at inference time.
The dominant RAG paradigm encodes a document corpus into dense vector
embeddings, retrieves the top-$k$ most similar passages to a query, and
concatenates them as context for the LLM. While effective for
fact-retrieval tasks, this flat chunk-based approach has two key
limitations in the academic setting:

\begin{enumerate}
  \item \textbf{Loss of relational structure.} Research papers contain
  rich typed relationships---a method proposes an algorithm, is evaluated
  on a dataset, reports a metric, and is compared to a baseline. Flat
  chunking severs these connections.
  \item \textbf{Poor cross-paper synthesis.} When a user asks ``which of
  these papers achieves the highest AP on COCO?'', independent text
  chunks do not naturally enable structured comparison.
\end{enumerate}

Graph RAG \cite{edge2024graphrag} partially addresses the first limitation
by traversing a knowledge graph, but pure graph traversal can miss
relevant contextual passages not captured in extracted entities.

We propose \textbf{Research-RAG}, a system unifying three retrieval
modalities: (1)~\emph{dense} semantic search via FAISS
\cite{johnson2019billion}, (2)~\emph{sparse} BM25
\cite{robertson2009probabilistic} for exact-term recall, and
(3)~\emph{graph-neighbourhood} retrieval from an LLM-extracted knowledge
graph. A linear weighted fusion combines dense and sparse scores; the
merged evidence is supplied to an LLM to produce a structured,
evidence-mapped JSON answer.

The main contributions of this paper are:
\begin{itemize}
  \item A four-phase pipeline (Parse $\to$ Extract $\to$ Build $\to$
  Index) that ingests arbitrary academic PDFs (\S\ref{sec:system}).
  \item A hybrid retrieval module that fuses dense, BM25, and
  graph-neighbourhood signals (\S\ref{sec:retrieval}).
  \item A multi-paper answer generator with grounding verification
  (\S\ref{sec:generation}).
  \item An empirical evaluation on real papers with an LLM-as-judge
  protocol, comparing all three systems on faithfulness and completeness
  (\S\ref{sec:experiments}).
\end{itemize}

% =========================================================
\section{Related Work}
\label{sec:related}

\subsection{Retrieval-Augmented Generation}
Lewis et al.~\cite{lewis2020rag} introduced RAG for open-domain question
answering. Subsequent work explored denser retrieval \cite{karpukhin2020dpr},
multi-hop reasoning \cite{xiong2021multihop}, and modular architectures
\cite{gao2023modular}.

\subsection{Graph-Enhanced RAG}
KGRAG \cite{hu2023kgrag} and GraphRAG \cite{edge2024graphrag} inject
knowledge graph triples into the retrieval context. Our work differs in
that we maintain both a full text index and a fine-grained, paper-specific
knowledge graph, enabling both passage-level and entity-level retrieval.

\subsection{Scientific Document Understanding}
SciREX \cite{jain2020scirex} and SPECTER \cite{cohan2020specter} address
structured information extraction and embedding for scientific papers.
Our extraction schema extends these with an evidence map linking each
extracted field to specific page numbers and text snippets.

% =========================================================
\section{System Architecture}
\label{sec:system}

The system consists of four sequential phases executed at indexing time
and a unified retrieval-generation pipeline invoked at query time.

\subsection{Phase 1: PDF Parsing (\texttt{parser.py})}
\label{sec:parsing}

Raw PDFs are processed using PyMuPDF (\texttt{fitz}) to extract text at
the block level, preserving page numbers, font sizes, and bold flags. A
heuristic identifies section headers: any text block whose font size
exceeds the modal body font size by more than 0.5~pt, or that uses a
bold typeface and contains fewer than fifteen words, is classified as a
header (Algorithm~\ref{alg:parse}).

\begin{algorithm}
\caption{Section-Aware Chunking}
\label{alg:parse}
\begin{algorithmic}[1]
\Require PDF path, chunk\_size $C$, overlap $O$
\Ensure List of \texttt{DocumentChunk} objects
\State body\_size $\leftarrow$ modal font size over first 3 pages
\State current\_section $\leftarrow$ ``Main'', buffer $\leftarrow$ ``'', pages $\leftarrow$ []
\For{each page block $b$}
  \If{is\_header($b$, body\_size)}
    \State current\_section $\leftarrow b$.text
  \EndIf
  \If{$|$buffer$|$ $>$ $C$ \textbf{or} (section changed and $|$buffer$|$ $>$ $C/2$)}
    \State emit chunk(buffer, pages, current\_section)
    \State buffer $\leftarrow$ last $O$ chars of buffer $+$ $b$.text
  \Else
    \State buffer $\leftarrow$ buffer $+$ $b$.text
  \EndIf
  \State pages.append($b$.page)
\EndFor
\State emit chunk(buffer, pages, current\_section)
\end{algorithmic}
\end{algorithm}

This produces \texttt{DocumentChunk} Pydantic objects carrying
\texttt{page\_start}, \texttt{page\_end}, \texttt{section}, and
\texttt{text}. Default configuration: chunk~size~1500~chars,
overlap~200~chars.

\subsection{Phase 2: Structured Entity Extraction (\texttt{multi\_extractor.py})}
\label{sec:extraction}

Chunks are priority-ranked by section keywords (abstract, method,
result, evaluation, limitation) and submitted to the LLM with a
structured extraction prompt. The extraction schema
(\texttt{schemas.py}) defines a \texttt{PaperExtraction}
Pydantic model comprising:

\begin{itemize}
  \item \textbf{Bibliographic metadata:} title, authors, venue, year, DOI.
  \item \textbf{Method block:} name, purpose, typed inputs/outputs, core
  idea, algorithm steps, assumptions, limitations, hyperparameters, and
  code/data links.
  \item \textbf{Evaluation block:} dataset names, metric names, and
  Result objects (dataset, metric, value, baseline, delta).
  \item \textbf{Provenance} (evidence map): page numbers and source
  snippets for each extracted field.
\end{itemize}

A \texttt{RobustStr} validator coerces unexpected JSON sub-objects,
making the schema resilient to LLM formatting variations. Two LLM
back-ends are supported---Groq (\texttt{llama-3.3-70b-versatile}) and
Gemini (\texttt{gemini-2.5-flash})---selectable via a \texttt{--backend}
flag.

\subsection{Phase 3: Knowledge Graph Construction (\texttt{graph\_builder.py})}
\label{sec:graph}

The structured extraction is translated into a directed typed
knowledge graph using NetworkX. Each entity becomes a typed node;
typed directed edges encode semantic relationships.
Table~\ref{tab:graph_schema} lists the node and edge types.

\begin{table}[htbp]
  \caption{Knowledge Graph Schema}
  \label{tab:graph_schema}
  \centering
  \begin{tabular}{llp{3.6cm}}
    \toprule
    \textbf{Type} & \textbf{Category} & \textbf{Description} \\
    \midrule
    Paper       & Node & Root node per paper \\
    Method      & Node & Proposed algorithm \\
    InputType   & Node & Typed model input \\
    OutputType  & Node & Typed model output \\
    Dataset     & Node & Evaluation dataset \\
    Metric      & Node & Evaluation metric \\
    Result      & Node & Numeric result entry \\
    Baseline    & Node & Compared baseline \\
    Limitation  & Node & Stated limitation \\
    Assumption  & Node & Stated assumption \\
    Resource    & Node & Code/data URL \\
    \midrule
    proposes      & Edge & Paper $\to$ Method \\
    takes\_input  & Edge & Method $\to$ InputType \\
    produces      & Edge & Method $\to$ OutputType \\
    evaluated\_on & Edge & Method $\to$ Dataset \\
    reports\_metric & Edge & Method $\to$ Metric \\
    reports       & Edge & Method $\to$ Result \\
    compared\_to  & Edge & Result $\to$ Baseline \\
    has           & Edge & Paper $\to$ Limitation/Assumption \\
    provides      & Edge & Paper $\to$ Resource \\
    \bottomrule
  \end{tabular}
\end{table}

Graph construction is additive: when multiple papers are loaded, their
subgraphs are merged into a single global directed graph, enabling
cross-paper entity queries via traversal.

\subsection{Phase 4: Hybrid Index Construction (\texttt{graph\_index.py})}
\label{sec:index}

The \texttt{HybridIndex} class constructs two complementary indices.

\subsubsection{Dense Chunk Index}
All chunk texts are encoded using the \texttt{all-MiniLM-L6-v2}
sentence-transformer \cite{reimers2019sbert} into 384-dimensional
unit-norm vectors, stored in a FAISS \texttt{IndexFlatIP}.

\subsubsection{Sparse BM25 Index}
The same chunks are tokenised and indexed with \texttt{BM25Okapi}
\cite{robertson2009probabilistic} for term-frequency and
inverse-document-frequency signals.

\subsubsection{Graph Node Index}
Each graph node is serialised as a textual description and encoded
with the same sentence transformer into a second FAISS index, enabling
semantic navigation of the graph.

% =========================================================
\section{Hybrid Retrieval}
\label{sec:retrieval}

At query time, the system executes three retrieval calls and merges the
evidence into a single context block.

\subsection{Dense Retrieval}
The query is encoded and cosine similarities are obtained via FAISS
inner product search over the chunk index, returning $3k$ candidates.

\subsection{Sparse BM25 Retrieval}
The query is tokenised and BM25 scores computed over the full corpus,
normalised to $[0, 1]$ against the top-scoring document.

\subsection{Hybrid Score Fusion}
Dense and sparse scores are linearly interpolated:
\begin{equation}
  s_\text{hybrid}(d, q) = \alpha \cdot s_\text{dense}(d, q)
                        + (1 - \alpha) \cdot s_\text{BM25}(d, q)
  \label{eq:fusion}
\end{equation}
where $\alpha = 0.5$ by default. The union of candidate sets is
re-ranked; the top $k$ are retained.

\subsection{Graph Neighbourhood Retrieval}
The query is encoded and top-$k$ matching graph nodes are identified
via FAISS search over the graph node index. For each matched node, the
1-hop ego-graph (undirected view) is extracted via NetworkX and
serialised as typed triples.

% =========================================================
\section{Answer Generation and Grounding}
\label{sec:generation}

\subsection{Multi-Paper Prompt Construction}
Retrieved text chunks and graph triples are concatenated with clear
section delimiters. The LLM is instructed to answer in strict JSON
containing: \texttt{reasoning\_summary}, \texttt{narrative\_answer},
\texttt{structured\_summary} (key findings, comparisons, agreements,
disagreements), \texttt{evidence} (page-level citations), and
\texttt{insufficiency} (flag + missing fields).

\subsection{Grounding Verification}
A second LLM call fact-checks the generated answer against the
retrieved context, returning a \texttt{grounding\_score} $\in [0, 1]$,
\texttt{hallucinated\_claims}, and \texttt{grounded\_claims}. This
two-pass design separates generation from verification.

% =========================================================
\section{Experimental Evaluation}
\label{sec:experiments}

\subsection{Setup}
\label{sec:setup}

We evaluate on \textbf{5 real PDFs} from the project's paper corpus,
covering few-shot object detection, incremental object detection, image
captioning, and benchmark comparison:

\begin{itemize}
  \item \textbf{P1} --- ``Meta Faster R-CNN'' \cite{han2021fasterrcnn}
  \item \textbf{P2} --- ``YOLO-IOD'' \cite{zhang2026yoloincrementalod}
  \item \textbf{P3} --- ``C4PS: Image Enhancement and Multilingual Captioning''
  \item \textbf{P4} --- ``Comparative Analysis of DL Image Detection'' \cite{srivastava2021comparative}
  \item \textbf{P5} --- Lab notebook (used only in indexing)
\end{itemize}

The PDF parser produced \textbf{114 chunks} in total. The knowledge graph
contained \textbf{58 nodes} and \textbf{53 edges} after merging all five
papers' extractions.

A set of \textbf{12 questions} was authored targeting the four main
papers, divided into:
\begin{itemize}
  \item \textbf{4 factual} (Q01--Q04): direct lookup of a specific fact
  (dataset, metric, architecture component).
  \item \textbf{4 reasoning} (Q05--Q08): requires interpreting or
  summarising implicit information (limitations, problem motivation).
  \item \textbf{4 comparison} (Q09--Q12): cross-paper analysis.
\end{itemize}

\subsection{Baselines}
\label{sec:baselines}

\begin{itemize}
  \item \textbf{Simple~RAG}: Dense-only retrieval (same
  \texttt{all-MiniLM-L6-v2} encoder and FAISS index, no BM25, no graph).
  \item \textbf{Graph~RAG}: Graph-neighbourhood retrieval only---FAISS
  search over the graph node index, returning ego-graph triples. No raw
  text chunks.
  \item \textbf{Hybrid~RAG}: The full Research-RAG system: dense + BM25
  fusion ($\alpha = 0.5$) + graph neighbourhood.
\end{itemize}

All three systems use the same generation LLM:
\texttt{llama-3.3-70b-versatile} (temperature~=~0) via Groq.
Retrieval context is capped at 8~000 characters for generation and
3~500 for judging.

\subsection{Evaluation Protocol (LLM-as-Judge)}
\label{sec:metrics}

Each system's answer is graded by an independent \texttt{llama-3.3-70b}
judge call with access only to the retrieved context (not the full PDF):

\begin{itemize}
  \item \textbf{Faithfulness} ($F$): fraction of claims in the answer
  traceable to the retrieved context ($=1.0$: all claims supported;
  $=0.0$: significant hallucination).
  \item \textbf{Completeness} ($C$): fraction of the question addressed
  ($=1.0$: fully answered; $=0.0$: not answered).
  \item \textbf{Combined} = $(F + C) / 2$.
\end{itemize}

\subsection{Results}
\label{sec:results}

\subsubsection{Overall Aggregate}

Table~\ref{tab:results_main} reports mean scores over all 12 questions.
Simple~RAG achieves the highest overall Combined score (0.8458), driven
by consistently high faithfulness (0.9833) and good completeness
(0.7083). Graph~RAG has faithfulness identical to Simple~RAG but suffers
significantly on completeness (0.4833), leading to a lower Combined
(0.7333). Hybrid~RAG sits in between (Combined~=~0.7958).

\begin{table}[htbp]
  \caption{Overall Aggregate Scores (12 questions, mean)}
  \label{tab:results_main}
  \centering
  \begin{tabular}{lccc}
    \toprule
    \textbf{System} & \textbf{Faithfulness} & \textbf{Completeness}
      & \textbf{Combined} \\
    \midrule
    Simple RAG  & 0.9833 & 0.7083 & 0.8458 \\
    Graph RAG   & 0.9833 & 0.4833 & 0.7333 \\
    Hybrid RAG  & 0.9083 & 0.6833 & 0.7958 \\
    \bottomrule
  \end{tabular}
\end{table}

\subsubsection{By Question Type}

Table~\ref{tab:results_breakdown} decomposes scores by question type.

\begin{table}[htbp]
  \caption{Combined Score by Question Type}
  \label{tab:results_breakdown}
  \centering
  \begin{tabular}{lccc}
    \toprule
    \textbf{Question Type} & \textbf{Simple RAG} & \textbf{Graph RAG}
      & \textbf{Hybrid RAG} \\
    \midrule
    Factual     & 0.8500 & 0.8000 & 0.7750 \\
    Reasoning   & 0.8750 & 0.6125 & 0.8500 \\
    Comparison  & 0.8125 & 0.7875 & 0.7625 \\
    \midrule
    \textbf{Overall} & \textbf{0.8458} & \textbf{0.7333} & \textbf{0.7958} \\
    \bottomrule
  \end{tabular}
\end{table}

Key observations:

\begin{itemize}
  \item \textbf{Reasoning questions}: Hybrid~RAG (0.8500) nearly matches
  Simple~RAG (0.8750) and \emph{substantially} outperforms Graph~RAG
  (0.6125, a gap of $+0.2375$). Pure graph retrieval fails here because
  reasoning questions require prose context that is not captured in
  structured triples. For example, on Q04 (``What image enhancement model
  is used in C4PS?''), Graph~RAG scored Completeness~=~0.00 because the
  graph had no \texttt{Model:Real-ESRGAN} node and could not answer from
  triples alone; Simple~RAG retrieved the relevant text chunk directly.
  Hybrid~RAG retrieved both, achieving Completeness~=~0.70.

  \item \textbf{Structured-fact questions}: Graph~RAG achieves
  Combined~=~1.00 on Q02 (``What datasets does YOLO-IOD use?'') and Q12
  (``What are the inputs/outputs of YOLO-IOD and Meta Faster R-CNN?''),
  where the exact answers were encoded as edges in the graph
  (\texttt{Method:YOLO-IOD --[evaluated\_on]--> Dataset:LoCo COCO}).
  Simple~RAG scored Combined~=~0.50 on Q12 because the prose chunks did
  not contain explicit input/output tables.

  \item \textbf{Comparison questions}: All three systems perform
  comparably (0.76--0.81), indicating that cross-paper synthesis is
  primarily a generation challenge rather than a retrieval one.
\end{itemize}

\subsubsection{Per-Question Detail}

Table~\ref{tab:results_perq} reports per-question faithfulness (F) and
completeness (C) for all three systems.

\begin{table}[htbp]
  \caption{Per-Question Scores (F = Faithfulness, C = Completeness)}
  \label{tab:results_perq}
  \scriptsize
  \centering
  \renewcommand{\arraystretch}{1.1}
  \begin{tabular}{l|cc|cc|cc}
    \toprule
    \textbf{Q} & \multicolumn{2}{c|}{\textbf{Simple RAG}}
               & \multicolumn{2}{c|}{\textbf{Graph RAG}}
               & \multicolumn{2}{c}{\textbf{Hybrid RAG}} \\
               & F & C & F & C & F & C \\
    \midrule
    Q01 [fact]  & 1.00 & 0.70 & 1.00 & 0.60 & 0.80 & 0.40 \\
    Q02 [fact]  & 1.00 & 0.80 & 1.00 & 1.00 & 1.00 & 0.80 \\
    Q03 [fact]  & 1.00 & 0.50 & 1.00 & 0.80 & 0.90 & 0.80 \\
    Q04 [fact]  & 1.00 & 0.80 & 1.00 & 0.00 & 0.80 & 0.70 \\
    \hline
    Q05 [reas]  & 1.00 & 1.00 & 1.00 & 0.50 & 1.00 & 0.80 \\
    Q06 [reas]  & 0.90 & 0.80 & 0.80 & 0.60 & 0.90 & 0.80 \\
    Q07 [reas]  & 1.00 & 0.80 & 1.00 & 0.00 & 0.90 & 0.70 \\
    Q08 [reas]  & 0.90 & 0.60 & 1.00 & 0.00 & 0.90 & 0.80 \\
    \hline
    Q09 [comp]  & 1.00 & 1.00 & 1.00 & 0.80 & 1.00 & 1.00 \\
    Q10 [comp]  & 1.00 & 0.50 & 1.00 & 0.50 & 0.80 & 0.60 \\
    Q11 [comp]  & 1.00 & 1.00 & 1.00 & 0.00 & 0.90 & 0.80 \\
    Q12 [comp]  & 1.00 & 0.00 & 1.00 & 1.00 & 1.00 & 0.00 \\
    \bottomrule
  \end{tabular}
\end{table}

\subsection{Error Analysis}
\label{sec:error}

\paragraph{Graph~RAG completeness failures}
On Q04, Q07, Q08, Q11 and Q12, Graph~RAG returns
Completeness~=~0.00. In all cases, the required information was present
in the source PDF as prose but was not captured in the knowledge graph
during extraction. The graph extractor relies on an LLM-driven
prompt that may miss fine-grained information (e.g., specific model
names in a table, or per-language translation details). This confirms
that pure graph retrieval is insufficient for questions that depend on
paragraph-level context.

\paragraph{Hybrid~RAG completeness failures}
Hybrid~RAG also scores Completeness~=~0.00 on Q12, the same structured
input/output question. Investigation shows that the BM25 component
retrieved chunks discussing YOLO's general architecture and COCO
benchmarks rather than the explicit input/output sections, which
appeared only in the artifact JSON (and thus in the graph). Increasing
the graph component's weight ($\alpha < 0.5$) would address this case.

\paragraph{Faithfulness is uniformly high}
All three systems maintain Faithfulness $\geq 0.80$ across all
questions. This confirms that the system's explicit instruction to
answer only from context is effective; the LLM rarely invents facts
beyond what is retrieved.

\subsection{Qualitative Example (Q08)}
\label{sec:qualitative}

Consider Q08: ``How does C4PS handle multilingual caption generation?''

\begin{itemize}
  \item \textbf{Graph~RAG}: Retrieved only the structured node
  \texttt{OutputType:Multilingual~Captions} and its edge back to
  \texttt{Method:C4PS}. The answer correctly stated that C4PS produces
  multilingual captions but could not explain the mechanism
  (Completeness~=~0.00).
  \item \textbf{Simple~RAG}: Retrieved a text chunk from the paper's
  system description but the chunk was cut before the MarianMT/NLLB
  details (Completeness~=~0.60).
  \item \textbf{Hybrid~RAG}: Retrieved the same text plus an additional
  chunk with the translation architecture details (MarianMT for French/
  Spanish/German, IndicTransWrapper for Indian languages), producing the
  most complete answer (Completeness~=~0.80).
\end{itemize}

This illustrates the complementary nature of the three modalities: the
graph confirms the existence of multilingual output, the dense component
finds the architectural prose, and BM25 boosts the chunk containing rare
technical terms like ``MarianMT'' and ``NLLB-200''.

% =========================================================
\section{Discussion}
\label{sec:discussion}

\subsection{When Does Hybrid RAG Win?}
The experiments show that Hybrid~RAG's advantage is most pronounced on
\textbf{reasoning questions} that require prose context not captured in
structured extraction. For \textbf{structured-fact questions} where the
exact answer is a typed graph edge, Graph~RAG is superior.
Simple~RAG retains relevance due to consistently high faithfulness and
good completeness on prose-heavy questions.

The key insight is that the three modalities are complementary rather
than competitive: dense retrieval finds semantically similar passages,
BM25 boosts rare technical terms, and graph traversal surfaces exact
structured facts. A dynamically weighted fusion that adjusts $\alpha$
based on query type could capture the best of all three.

\subsection{Limitations}
\begin{itemize}
  \item \textbf{Extraction dependence.} Graph quality is bounded by LLM
  extraction accuracy. On Q04, the extractor did not create a node for
  \texttt{Model:Real-ESRGAN}, so graph retrieval failed completely.
  \item \textbf{Judge-as-Oracle.} The LLM judge grades answers against
  the \emph{retrieved context}, not the full PDF. A question that
  Simple~RAG answers correctly because the right chunk was retrieved will
  score high even if the answer omits facts present elsewhere in the PDF.
  Human evaluation against the full PDF would provide stronger ground truth.
  \item \textbf{Scale.} Results are from a 12-question benchmark over 4
  papers. Larger-scale evaluation is needed.
\end{itemize}

\subsection{Future Work}

\begin{enumerate}
  \item \textbf{Adaptive fusion weights}: learn $\alpha$ per query type
  (structured-fact queries $\to$ lower $\alpha$; prose reasoning $\to$
  higher $\alpha$).
  \item \textbf{Cross-paper entity resolution}: merge equivalent nodes
  (e.g., two papers' \texttt{Dataset:COCO} nodes) across papers.
  \item \textbf{Incremental graph updates}: add papers without full
  re-indexing.
  \item \textbf{Human evaluation}: expert annotators scoring answers
  against the original PDF.
\end{enumerate}

% =========================================================
\section{Conclusion}
\label{sec:conclusion}

We presented Research-RAG, a hybrid retrieval-augmented generation
system for academic question answering over multiple research papers. By
fusing dense semantic search, BM25 sparse retrieval, and knowledge-graph
neighbourhood traversal, the system combines complementary strengths.
Real experiments with 12 questions across 4 papers and LLM-as-judge
evaluation reveal that: (1)~Simple~RAG achieves the highest overall
Combined score (0.8458) due to consistently complete prose retrieval; (2)~
Graph~RAG excels on structured-fact questions where answers are encoded
as typed edges (Combined~=~1.00 on Q02 and Q12) but fails on
reasoning queries (Combined~=~0.6125, $-0.2375$ vs.\ Hybrid~RAG); (3)~
Hybrid~RAG matches Simple~RAG on reasoning (0.8500 vs.\ 0.8750) while
recovering graph-RAG's failures on prose-heavy questions.
The results confirm the complementary nature of the three modalities and
motivate adaptive fusion strategies as future work.

% =========================================================
\section*{Acknowledgments}
The authors thank the open-source communities behind PyMuPDF,
FAISS, NetworkX, sentence-transformers, rank-bm25, and Groq for
the foundational libraries upon which this system is built.

% =========================================================
\begin{thebibliography}{99}

\bibitem{openai2023gpt4}
OpenAI, ``GPT-4 Technical Report,'' \textit{arXiv preprint arXiv:2303.08774}, 2023.

\bibitem{team2023gemini}
Gemini Team, Google, ``Gemini: A Family of Highly Capable Multimodal
Models,'' \textit{arXiv preprint arXiv:2312.11805}, 2023.

\bibitem{ji2023survey}
Z.~Ji, N.~Lee, R.~Frieske, T.~Yu, D.~Su, Y.~Xu, E.~Ishii, Y.~Bang,
A.~Madotto, and P.~Fung, ``Survey of Hallucination in Natural Language
Generation,'' \textit{ACM Comput. Surv.}, vol.~55, no.~12, pp.~1--38, 2023.

\bibitem{lewis2020rag}
P.~Lewis et al., ``Retrieval-Augmented Generation for Knowledge-Intensive
NLP Tasks,'' in \textit{Proc. NeurIPS}, 2020, pp.~9459--9474.

\bibitem{karpukhin2020dpr}
V.~Karpukhin et al., ``Dense Passage Retrieval for Open-Domain Question
Answering,'' in \textit{Proc. EMNLP}, 2020, pp.~6769--6781.

\bibitem{xiong2021multihop}
W.~Xiong et al., ``Answering Complex Open-Domain Questions with
Multi-Hop Dense Retrieval,'' in \textit{Proc. ICLR}, 2021.

\bibitem{gao2023modular}
Y.~Gao et al., ``Retrieval-Augmented Generation for Large Language
Models: A Survey,'' \textit{arXiv preprint arXiv:2312.10997}, 2023.

\bibitem{edge2024graphrag}
D.~Edge et al., ``From Local to Global: A Graph RAG Approach to
Query-Focused Summarization,'' \textit{arXiv preprint arXiv:2404.16130}, 2024.

\bibitem{hu2023kgrag}
X.~Hu et al., ``Think-on-Graph: Deep and Responsible Reasoning of Large
Language Models on Knowledge Graphs,''
\textit{arXiv preprint arXiv:2307.07697}, 2023.

\bibitem{jain2020scirex}
S.~Jain et al., ``SciREX: A Challenge Dataset for Document-Level
Information Extraction,'' in \textit{Proc. ACL}, 2020, pp.~7506--7516.

\bibitem{cohan2020specter}
A.~Cohan et al., ``SPECTER: Document-level Representation Learning
using Citation-informed Transformers,'' in \textit{Proc. ACL}, 2020,
pp.~2270--2282.

\bibitem{johnson2019billion}
J.~Johnson, M.~Douze, and H.~J{\'{e}}gou, ``Billion-Scale Similarity
Search with GPUs,'' \textit{IEEE Trans. Big Data}, vol.~7, no.~3,
pp.~535--547, 2021.

\bibitem{robertson2009probabilistic}
S.~Robertson and H.~Zaragoza, ``The Probabilistic Relevance Framework:
BM25 and Beyond,'' \textit{Found. Trends Inf. Retr.}, vol.~3, no.~4,
pp.~333--389, 2009.

\bibitem{reimers2019sbert}
N.~Reimers and I.~Gurevych, ``Sentence-BERT: Sentence Embeddings using
Siamese BERT-Networks,'' in \textit{Proc. EMNLP}, 2019, pp.~3982--3992.

\bibitem{han2021fasterrcnn}
G.~Han, S.~Huang, J.~Ma, Y.~He, and S.-F.~Chang, ``Meta Faster R-CNN:
Towards Accurate Few-Shot Object Detection with Attentive Feature
Alignment,'' \textit{arXiv preprint arXiv:2104.07719}, 2021.

\bibitem{zhang2026yoloincrementalod}
S.~Zhang et al., ``YOLO-IOD: Towards Real Time Incremental Object
Detection,'' \textit{arXiv preprint arXiv:2512.22973}, 2025.

\bibitem{srivastava2021comparative}
S.~Srivastava et al., ``Comparative analysis of deep learning image
detection algorithms,'' \textit{J. Big Data}, vol.~8, no.~66, 2021.

\end{thebibliography}

\end{document}
